% Solution: Milgrom & Roberts (1992) Ch. 8 Food for Thought 3 (auditors). Input by \solitem. Check: check_T1.py (item 6).
\begin{exsol}{Milgrom \& Roberts Ch.~8, Food for Thought 3}{284}{Foundations}{2}
\asked Why do large corporations use big audit firms with many clients, and avoid a firm for which one
client would be too large a share of its business?

\method The chapter's reputation model (pp.~259--268; efficiency-wage condition, pp.~250--255).
A party who can cheat for a one-time gain $G$ stays honest iff the value of the reputation he would lose
exceeds $G$: in the efficiency-wage form, $p\,(w - \bar w)\,N \ge G$. We read the auditor's reputation as the
trust-based rent $v$ that investors' willingness to rely on its reports earns it. The answer is an
analytical argument, not a computation; the chapter prints no answer to this question.

\solution
\textbf{1. What auditors sell.} Investors cannot verify the books themselves. An audit is valuable only
if the auditor reports honestly even though management picks and pays him. This is the trust game of
Ch.~8: investors offer trust, and the auditor can fail to honour it by shading the report to please
management (gain $G$). Contracts cannot specify honest reporting in detail, so reputation has to enforce it
(``Reputations as contract enforcers'', p.~259).

\textbf{2. Many clients make honesty cheaper to enforce.}
\begin{itemize}
  \item \emph{Reputation is firm-wide and indivisible.} A shaded report for one client, once detected, damages
  the auditor's credibility with all clients and with investors, so the multiplier $N$ applies to the whole
  stream of rents (compare the Maghribi traders, who punished cheaters group-wide, p.~252). With $k$ clients
  paying a rent $v$ each, honesty requires $p\,k\,v\,N \ge G_i$ for every client $i$. The left side grows
  with $k$, while the temptation $G_i$ from any one client is only part of the firm's income.
  \item \emph{A dominant client changes both sides.} If one client is a large share of revenue, its
  favour is worth more ($G_i$ large), and the auditor has more to lose from displeasing that client
  than from being caught. The client can credibly threaten to leave and take the auditor's income with it:
  a firm that depends on one buyer faces the buyer's hold-up. With many clients, no single threat is
  decisive.
  \item \emph{Reputation outlives individuals.} The brand passes from one partner to the next
  (Tadelis \S10.5.3, pp.~207--208), and large firms monitor and discipline their own partners, which keeps
  the detection probability $p$ high.
\end{itemize}

\textbf{3. Why management goes along.} Managers want the credibility of a clean opinion in order to raise
capital, so they choose a firm whose reputation investors recognise. A small firm dependent on one client
has no reputation to lose, so investors would discount its opinion and the audit would be worthless to
the very manager who buys it.

\textbf{4. Auxiliary reasons.} Economies of scale, specialised knowledge, pooling of litigation risk and
the capacity to audit a large corporation's worldwide records also favour big firms. They explain size,
but not the reluctance to hire a firm with a large client share, which is the puzzle.

\interp Audit independence is bought by making the auditor's stake in its whole reputation large
relative to the gain from any single favour. This is also why regulators limit the share of a firm's fees
that may come from one audit client.

\verify \texttt{check\_T1.py} (item~6) illustrates the inequality $p\,k\,v\,N \ge G$ with assumed numbers
($G = 10$, $p = \frac12$, $v = 1$, $N = 6$): honesty needs $k \ge 4$ clients, and $k = 3$ fails. The numbers
are illustrative, not from the book.

\answer{The auditor's reputation for honest reporting is a firm-wide asset. A detected shaded report loses
all clients' rents, while the gain from favouring one client is limited. A firm dominated by one client has
little to lose from displeasing it and much to lose by losing it, so its independence is weak and investors
discount its opinion. Large audit firms with many clients make honesty incentive compatible, and
management gains from the credibility.}
\end{exsol}
